% --------------------------------------------------
% 基本数学与环境
% --------------------------------------------------
\usepackage{amsmath,amssymb,amsfonts,amsthm,mathtools,bm}
\usepackage{array,booktabs,longtable,multicol,multirow,makecell,tabularx,threeparttable}
\usepackage{enumitem}
\setlist[enumerate]{noitemsep, topsep=2pt}
\setlist[itemize]{noitemsep, topsep=2pt}

% --------------------------------------------------
% 版面与引用
% --------------------------------------------------
\usepackage[margin=0.9in]{geometry}
\usepackage[numbers,sort&compress]{natbib}
\usepackage[colorlinks,citecolor=purple,linkcolor=blue]{hyperref}
\usepackage{cleveref}
\usepackage{float}
\usepackage{caption,subcaption}
\usepackage{appendix}
\usepackage{authblk}
\usepackage{parskip}
\renewcommand{\baselinestretch}{1.2}

% --------------------------------------------------
% 算法环境
% --------------------------------------------------
% \usepackage[linesnumbered,boxruled,lined,commentsnumbered]{algorithm2e} % 如果只用 algorithmic，可切换掉这一行
\usepackage{algorithm,algorithmic}
\usepackage{eqparbox}
\renewcommand\algorithmiccomment[1]{\hfill\texttt{\#}\ \eqparbox{COMMENT}{\texttt{#1}}}

% 自定义 subroutine 浮动体
\usepackage{newfloat}
\newfloat{subroutine}{htbp}{los}
\floatname{subroutine}{Subroutine}
\crefname{subroutine}{subroutine}{subroutines}
\Crefname{subroutine}{Subroutine}{Subroutines}
\crefname{line}{line}{lines}
\def\algorithmautorefname{Alg.}
\Crefname{algorithm}{Alg.}{Alg.}

% --------------------------------------------------
% Theorem 环境定义
% --------------------------------------------------
\newtheorem{theorem}{Theorem}[section]
\newtheorem{lemma}[theorem]{Lemma}
\newtheorem{proposition}[theorem]{Proposition}
\newtheorem{corollary}[theorem]{Corollary}
\newtheorem{assumption}[theorem]{Assumption}
\newtheorem{condition}[theorem]{Condition}
\newtheorem{definition}[theorem]{Definition}
\newtheorem{remark}[theorem]{Remark}
\newtheorem{example}[theorem]{Example}
\newtheorem{strategy}[theorem]{Strategy}
\newtheorem{question}[theorem]{Question}
\numberwithin{equation}{section}

% cleveref 自动命名
\def\subsectionautorefname{Section}
\def\definitionautorefname{Definition}
\def\corollaryautorefname{Corollary}
\def\lemmaautorefname{Lemma}
\def\assumptionautorefname{Assumption}
\def\propositionautorefname{Proposition}
\def\tableautorefname{Table}
\def\strategyautorefname{Strategy}
\def\remarkautorefname{Remark}
\def\conditionautorefname{Condition}

% --------------------------------------------------
% TikZ & PGFPlots
% --------------------------------------------------
\usepackage{tikz,pgfplots}
\usetikzlibrary{arrows.meta,calc,backgrounds,positioning}
\usepgfplotslibrary{fillbetween,patchplots}
\pgfplotsset{compat=newest}

% 自定义颜色与样式
\definecolor{ao(english)}{rgb}{0.0,0.5,0.0}
\definecolor{cadmiumgreen}{rgb}{0.0,0.42,0.24}
\definecolor{darkpastelgreen}{rgb}{0.01,0.75,0.24}

\tikzset{
  materia/.style={draw, fill=blue!20, text width=6em, text centered, minimum height=1.5em, drop shadow},
  etape/.style={materia, text width=8em, minimum width=10em, minimum height=3em, rounded corners, drop shadow},
  linepart/.style={draw, thick, color=black!50, -Latex, dashed},
  line/.style={draw, thick, color=black!50, -Latex},
  back group/.style={fill=white!20,rounded corners, draw=black!50, dashed, inner xsep=15pt, inner ysep=10pt},
}
\newcommand{\etape}[2]{node (p#1) [etape] {#2}}
\newcommand{\transreceptor}[3]{\path [linepart] (#1.east) -- node [above] {\scriptsize #2} (#3);}

% --------------------------------------------------
% 其他常用命令
% --------------------------------------------------
\newcommand{\xk}{x_k}
\newcommand{\Hk}{H_k}
\newcommand{\gk}{g_k}
\newcommand{\dk}{d_k}
\newcommand{\xkn}{x_{k+1}}
\newcommand{\gkn}{g_{k+1}}
\newcommand{\dkn}{d_{k+1}}
\newcommand{\cmark}{\ding{51}}
\newcommand{\xmark}{\ding{55}}

% 算法名缩写
\newcommand{\arc}{\texttt{CubicReg}}
\newcommand{\arca}{\texttt{CubicReg-A}}
\newcommand{\utr}{\texttt{UTR}}
\newcommand{\utrhvp}{\texttt{UTR-HVP}}
\newcommand{\iutr}{\texttt{iUTR(Hessian)}}
\newcommand{\iutrhvp}{\texttt{iUTR}}
\newcommand{\drsom}{\texttt{DRSOM}}
\newcommand{\drsomh}{\texttt{DRSOM-H}}
\newcommand{\hsodm}{\texttt{HSODM}}
\newcommand{\hsodmarc}{\texttt{HSODMArC}}
\newcommand{\hsodmhvp}{\texttt{HSODM-HVP}}
\newcommand{\atr}{\Cref{alg.accelerated utr}}
\newcommand{\atrms}{\Cref{alg.ms accelerated utr}}
\newcommand{\tr}{\textrm{TR}}
\newcommand{\trp}{\textrm{TR}_+}
\numberwithin{equation}{section}